% 05-experiences-methode (en): the protocol of the experiments, as the code runs it.
\chapter{The experiments: protocol and methods}\label{chap:05-experiences-methode}

The survey of \cref{chap:03-sondage-methode,chap:04-sondage-resultats} measures what papers publish. The
experiments put the flaws of 2018 to the test on real training runs, made for this report. They return to the
ground of \textcite{henderson2018matters}: simulated bodies in the MuJoCo physics simulator, bodies that
reinforcement learning algorithms learn to move forward. They do so with the algorithms of 2026 and with the
rules Joelle Pineau asked for: enough random seeds, confidence intervals, baselines tuned with the same care,
evaluation held out, code and data published.

This chapter fixes everything that comes before a result: the tools and their versions, the programs and their
settings, the tasks, the budgets, the random seeds, the evaluation protocol, the file format, the statistics
and the design of each experiment. It describes what the code does, line by line where that helps. It also
flags every place where the code departs from the written protocol. The results come in
\cref{chap:06-experiences-resultats}. \Cref{chap:08-logiciel} draws the software chain that links the runs to
the documents.

\section{A protocol written before the first training run}\label{sec:protocole}

The protocol of the experiments is the file \file{sim/SIM\_PROTOCOL.md}. Its first version is dated
8~November~2026. It was written before any training. It fixes the questions, the algorithms, the tasks, the
budgets, the number of random seeds, the evaluation and the analyses. It also lays down three rules of
conduct:
\begin{itemize}
  \item the random seeds are fixed before launching, from 0 to 29;
  \item no run is removed or relaunched because of its score;
  \item a run that crashes is relaunched once with the same random seed, then the incident is recorded.
\end{itemize}

Every later change enters a numbered, dated list of \enquote{deviations}, together with a note of what we had
already seen when we decided it. The protocol counts twenty of them as of 21~November~2026.
\Cref{chap:09-ecarts} gives all of them, with those of the survey and of the language-model study. This chapter
cites a deviation by its number when it changes a method. \Cref{sec:avant-apres} sorts the decisions by
whether they precede or follow the results they affect.

Throughout the report, a \emph{run} is one complete training of one program on one task with a given random
seed. The \emph{main} results are the ones the protocol announced: ten random seeds per pair of program and
task. The additions of deviation~12 (twenty random seeds, equal budget, replication) are given next to them.
Every divergence is stated.

\section{Tools and versions}\label{sec:outils}

\paragraph{Four frozen Python environments}
Four Python~3.11.18 environments were used, all on a Mac with an Apple M2 Max chip (macOS~27.2, arm64). Each
is frozen by a \texttt{pip freeze} taken on 15~November~2026, in \file{envs/} (\cref{tab:env},
\cref{chap:D-environnements}):
\begin{itemize}
  \item \file{sim/env} serves training and re-evaluation (\file{sim/run.py}, \file{sim/dose.py}). It holds
    Stable-Baselines3~2.9.2 \parencite{raffin2021sb3}, sb3-contrib~2.9.2, SBX~0.28.0 \parencite{raffin2023sbx}, JAX~0.10.2
    \parencite{bradbury2018jax}, Flax~0.12.8, Optax~0.2.8, PyTorch~2.15.0 \parencite{paszke2019pytorch},
    Gymnasium~1.1.1 \parencite{towers2024gymnasium}, MuJoCo~4.0.0 \parencite{todorov2012mujoco};
  \item \file{sim/env-an} computes every analysis: NumPy~2.4.6, SciPy~1.17.1, statsmodels~0.15.0,
    rliable~1.2.0 \parencite{agarwal2021precipice};
  \item \file{sim/env-zoo} runs the official script of RL Baselines3 Zoo~2.9.2 \parencite{raffin2020zoo}
    for experiment E6, with the same libraries as \file{sim/env};
  \item \file{sim/env-llm} serves the language-model study (\cref{chap:07-modeles-de-langue}).
\end{itemize}

\paragraph{What each tool contributes}
Gymnasium, the successor of OpenAI Gym maintained by the Farama Foundation, defines the interface of the
tasks: \texttt{reset} returns a first observation, \texttt{step} applies an action and returns the next
observation, the reward and two end-of-episode flags. MuJoCo computes the physics. It has been free and open
since 2022, whereas it required a paid licence in 2018. Stable-Baselines3 provides PPO written in PyTorch.
SBX (\enquote{Stable-Baselines Jax}) rewrites PPO, SAC and CrossQ in JAX with the same interface. sb3-contrib
provides TRPO. rliable computes interquartile means, their bootstrap intervals and probabilities of
improvement.

\paragraph{A version deviation}
The protocol announced Stable-Baselines3~2.6.0. Every training run used version~2.9.2, installed on
8~November at 10:11, before the first training run (deviation~14, item~c). The versions in
\file{envs/train.txt} are the ones that actually ran.

\paragraph{Reproducing on another machine}
On a given machine, a run is determined by its random seed: each training runs on a single thread
(\texttt{OMP\_NUM\_THREADS=1}, \texttt{torch.set\_num\_threads(1)}, single-thread XLA options for JAX). On
another processor, MuJoCo, JAX and PyTorch round differently, so the same random seed gives a different curve
there. A rerun elsewhere tests the statistics of the report. It recovers no curve identically.

\section{The tasks}\label{sec:taches}

The four tasks are those of \textcite{henderson2018matters}, in their current Gymnasium version:
\task{HalfCheetah-v5}, \task{Hopper-v5}, \task{Walker2d-v5} and \task{Swimmer-v5}. In each, MuJoCo simulates
an articulated body. At every step the agent observes a vector of positions and velocities, then chooses a
continuous action: the torque applied at each joint. The reward pays mostly for forward velocity, minus a cost
that grows with motor effort. For \task{Hopper} and \task{Walker2d}, it adds a bonus at every step where the
body stays upright. An episode stops after \num{1000}~steps at most.
\begin{itemize}
  \item \task{HalfCheetah} is a planar cheetah with six joints. It never falls, so its episodes always last
    \num{1000}~steps.
  \item \task{Hopper} is a one-legged hopper. The episode stops when it falls, which makes scores very uneven
    from one run to another.
  \item \task{Walker2d} is a planar two-legged walker, also stopped when it falls.
  \item \task{Swimmer} is a three-segment swimmer in a viscous fluid. Its scores are far lower than those of
    the other three tasks.
\end{itemize}

The \enquote{score} of an episode is its return: the sum of the rewards received from its first step to its
last. Scores from different tasks are not comparable. Every comparison in this report is made within one
task.

Experiment E6 also uses \task{Hopper-v4}, the previous version. The PPO settings that RL Baselines3 Zoo
publishes for Hopper are filed under this task name in the collection's \file{hyperparams/ppo.yml}: they were
tuned on that version.

\section{The programs}\label{sec:programmes}

Our experiments compare four \emph{programs} (four ready-to-use code bases) that implement three
\emph{algorithms}. A fifth program, TRPO, serves in one experiment only (B1). \Cref{tab:hyperparametres} gives
their settings.

\paragraph{PPO, twice}
PPO (\emph{Proximal Policy Optimization}, \cite{schulman2017ppo}) was the reference algorithm of 2018. It is
still a common point of comparison in 2026. It is also the direct ancestor of GRPO, the algorithm of
the language-model study. PPO is an \enquote{on-policy} algorithm: it collects a batch of transitions with the
current policy, uses it for a few gradient passes, then discards it. Let
$r_t(\theta) = \pi_\theta(a_t\mid s_t)/\pi_{\theta_{\text{old}}}(a_t\mid s_t)$ be the ratio of the
probabilities of the action under the new and the old policy. Let $\hat A_t$ be the advantage estimated by GAE
(\emph{generalized advantage estimation}): how much more the action taken returned than the value function
predicted for that state. PPO maximises the clipped objective
\begin{equation}\label{eq:ppo}
  L^{\text{CLIP}}(\theta) = \hat{\mathbb{E}}_t\Bigl[\min\bigl(r_t(\theta)\,\hat A_t,\;
  \operatorname{clip}(r_t(\theta), 1-\epsilon, 1+\epsilon)\,\hat A_t\bigr)\Bigr],
\end{equation}
with $\epsilon = 0.2$, plus a value-function error term. Two programs implement PPO:
\begin{itemize}
  \item \algo{PPO PyTorch}, the \texttt{PPO} class of Stable-Baselines3 (\texttt{ppo\_sb3} in the code);
  \item \algo{PPO JAX}, the \texttt{PPO} class of SBX (\texttt{ppo\_sbx} in the code).
\end{itemize}
Both share the same default values for every hyperparameter of the algorithm. They differ in weight
initialisation: Stable-Baselines3 initialises its layers with an orthogonal matrix
(\texttt{ortho\_init=True}), while SBX keeps the default initialisation of Flax (\texttt{ortho\_init=False}).
They also differ in the numerical library, its random generators and its rounding. Experiment E3 measures what
these differences cost.

\paragraph{SAC}
SAC (\emph{Soft Actor-Critic}, \cite{haarnoja2018sac}) has been the \enquote{off-policy} standard of continuous
control since 2019. It keeps every transition it has seen in a replay buffer and learns at every step from a
batch drawn from that buffer. It maximises the return augmented with the entropy of the policy,
$\sum_t \mathbb{E}\bigl[r_t + \alpha\,\mathcal{H}(\pi(\cdot\mid s_t))\bigr]$, with a coefficient $\alpha$
tuned automatically towards a target entropy. Two critic networks estimate the value of actions. Slow copies of
these networks, called target networks, stabilise learning. The program is the \texttt{SAC} class of SBX
(\texttt{sac}).

\paragraph{CrossQ}
CrossQ \parencite{bhatt2024crossq}, published at ICLR~2024, targets data efficiency, the score reached per step
of interaction with the environment. A descendant of SAC, it removes the target networks and applies batch normalisation in the critic
and the actor. It also widens the critic. The program is the \texttt{CrossQ} class of SBX (\texttt{crossq}).
Its critic has two layers of \num{1024}~units, instead of the paper's \num{2048}: SBX makes this choice to
compute faster. We kept it.

\paragraph{TRPO, for B1 only}
TRPO (\emph{Trust Region Policy Optimization}, \cite{schulman2015trpo}) is the algorithm of figure~5 of
\textcite{henderson2018matters}, the figure whose $t$ of 9 experiment B1 tries to explain
(\cref{sec:plan-b1}). It maximises the advantage weighted by the probability ratio under a constraint on the
mean Kullback-Leibler divergence between the old and the new policy,
$\hat{\mathbb{E}}_t\bigl[\mathrm{KL}(\pi_{\theta_{\text{old}}}\Vert\pi_\theta)\bigr] \le \delta$, solved by
conjugate gradient followed by a line search. The program is the \texttt{TRPO} class of sb3-contrib
(\texttt{trpo}), with all its default values.

\paragraph{TD3, planned then dropped}
The protocol planned TD3 if the budget allowed. The budget did not allow it: TD3 was not launched
(deviation~1).

\paragraph{Why the default values}
The four programs of the main grid keep the default values of their library. None was tuned. This choice
applies Pineau's request on baselines to the letter (flaw F7, baselines tuned less carefully than the method):
each program receives the same tuning effort, here none. The choice has a cost. A program well served by its
default values is favoured. The ranking of E2 holds for these settings. Experiment E6 measures, on one
case, what a published tuning changes. The only command that builds a program of the grid, in
\file{sim/run.py}, is \texttt{cls("MlpPolicy", TASK, seed=SEED, verbose=0)}, with no other argument.

\begin{table}[tbp]
\centering\small
\caption{The settings of the five programs, read in the code of the installed libraries
(Stable-Baselines3~2.9.2, sb3-contrib~2.9.2, SBX~0.28.0). They are the default values: no program was tuned.
An empty cell means the setting does not exist for that algorithm. The networks are multilayer perceptrons,
given by the widths of their hidden layers.}
\label{tab:hyperparametres}
\begin{tabularx}{\linewidth}{@{}>{\raggedright\arraybackslash}X
  >{\raggedright\arraybackslash}p{1.75cm}>{\raggedright\arraybackslash}p{1.75cm}
  >{\raggedright\arraybackslash}p{1.75cm}>{\raggedright\arraybackslash}p{1.95cm}
  >{\raggedright\arraybackslash}p{1.75cm}@{}}
\toprule
Setting & PPO PyTorch & PPO JAX & SAC & CrossQ & TRPO \\
\midrule
Learning rate & \num{3e-4} & \num{3e-4} & \num{3e-4} & \num{1e-3} & \num{1e-3} (critic) \\
Discount factor $\gamma$ & 0.99 & 0.99 & 0.99 & 0.99 & 0.99 \\
Steps collected per update & \num{2048} & \num{2048} & 1 & 1 & \num{2048} \\
Batch size & 64 & 64 & 256 & 256 & 128 \\
Passes over each batch of transitions & 10 & 10 & & & \\
GAE parameter $\lambda$ & 0.95 & 0.95 & & & 0.95 \\
Clipping $\epsilon$ & 0.2 & 0.2 & & & \\
Divergence constraint $\delta$ & & & & & 0.01 \\
Entropy coefficient & 0 & 0 & automatic & automatic & \\
Replay buffer & & & \num{1000000} & \num{1000000} & \\
Steps before learning starts & & & 100 & 100 & \\
Target network update $\tau$ & & & 0.005 & no target network & \\
Critic updates per actor update & & & 1 & 3 & 10 \\
Policy network & 64, 64, tanh & 64, 64, tanh & 256, 256, ReLU & 256, 256, ReLU & 64, 64, tanh \\
Value or critic network & 64, 64, tanh & 64, 64, tanh & 2 times 256, 256 & 2 times \num{1024}, \num{1024}, batch-normalised & 64, 64, tanh \\
Orthogonal initialisation & yes & no & & & yes \\
\bottomrule
\end{tabularx}
\end{table}

\paragraph{The published Zoo settings, for E6}
Experiment E6 pits PPO PyTorch with its default values against PPO PyTorch with the settings RL Baselines3
Zoo publishes for \task{Hopper-v4}. \file{sim/run.py} copies them into the constant
\texttt{ZOO\_PPO\_HOPPER}, identical to those of \file{rl\_zoo3/hyperparams/ppo.yml} in version~2.9.2
(deviation~5):
\begin{itemize}
  \item batches of 512~steps, minibatches of 32, 5~passes, $\gamma = 0.999$, $\lambda = 0.99$ for GAE;
  \item learning rate \num{9.80828e-5}, entropy coefficient \num{0.00229519}, value coefficient
    \num{0.835671}, gradient norm clipped at 0.7, clipping 0.2;
  \item networks of two layers of 256~units for the policy and the value, ReLU activation, no orthogonal
    initialisation, initial log standard deviation of the actions at $-2$;
  \item normalisation of observations and rewards by a running mean and standard deviation
    (\texttt{VecNormalize}).
\end{itemize}
Reward normalisation divides each reward by the standard deviation of a running discounted return. In the Zoo
recipe, the discount factor of that return is the algorithm's, 0.999. Our first pass
(\texttt{ppo\_sb3\_zoo}, option \texttt{-{}-zoo}) left the default value of \texttt{VecNormalize}, 0.99. The
second pass (\texttt{ppo\_sb3\_zoofix}, option \texttt{-{}-zoofix}) passes 0.999 (deviation~2). Both are kept.

\section{The budgets}\label{sec:budgets}

\paragraph{What the protocol fixes}
PPO PyTorch, PPO JAX and TRPO train for one million steps. SAC and CrossQ train for \num{200000}~steps. A step
is one transition of the environment: an observation, an action, a reward.

\paragraph{Why two budgets}
The two families of algorithms are not compared in the same regime. An on-policy algorithm discards its data
after each update, so it needs many steps. One million is the 2018 order of magnitude for these tasks. An
off-policy algorithm reuses its data at every step: papers from 2024 to 2026 compare SAC and CrossQ in a
\enquote{low-data} regime. The protocol set \num{300000}~steps. Deviation~1 brought them down to
\num{200000}, before the first training run, to keep the compute within a shared machine: we then measured
19~ms per step for these two algorithms.

\paragraph{The exact number of steps}
\file{sim/run.py} trains in blocks of \num{50000}~steps and evaluates after each block. The call is
\texttt{model.learn(50\_000, reset\_num\_timesteps=False)}, repeated until the budget is reached. PPO and TRPO
collect their transitions in batches of \num{2048}~steps and stop only at the end of a batch, so each block
lasts \num{51200}~steps. The saved models of PPO PyTorch, PPO JAX and TRPO count \num{1024000}~steps, while
the \texttt{step} column of the results reads \num{1000000}. The E6 models tuned like the Zoo, with batches of
512~steps, count \num{1003520}. SAC and CrossQ stop at \num{200000}~steps, as announced (deviation~14, item~d).

\paragraph{The equal budget}
The comparison of E2 is made at unequal budgets, which the report states each time. Deviation~12, item~iii,
adds a reading at equal budget: PPO is read at \num{200000}~steps on its evaluation curve, at the same point as
SAC and CrossQ. This rule was written before any PPO score at \num{200000}~steps had been looked at.

\section{The random seeds}\label{sec:graines}

\paragraph{What a random seed fixes}
A random seed is the starting number from which the computer draws all the randomness of a training run: the
initialisation of the networks, the initial positions of the simulated body, the noise of explored actions,
the batches drawn from the replay buffer. \file{sim/run.py} passes it to the constructor of the program
(\texttt{seed=SEED}), which derives from it the generators of NumPy, PyTorch or JAX and that of the training
environment. Same program, same random seed, same machine: same run.

\paragraph{Training random seeds}
The training random seeds are the integers 0 to 29, fixed before the first training run:
\begin{itemize}
  \item 0 to 9 for every pair of program and task, the main results;
  \item 10 to 19 for the consolidation of deviation~12, the twenty random seeds;
  \item 0 to 29 for PPO PyTorch on \task{HalfCheetah} (E1) and for TRPO on \task{HalfCheetah} (B1).
\end{itemize}
Random seeds 10 to 19 of PPO PyTorch on \task{HalfCheetah}, trained for E1 with the same code, serve as they
are in the consolidation.

\paragraph{Evaluation random seeds}
Evaluation uses random seeds \num{10000} to \num{10009}, one per evaluation episode. Training never uses them:
it uses random seeds 0 to 29. The same ten random seeds serve at every evaluation point, in every run and in
every physics, so two policies are always evaluated from the same ten starting positions.

\paragraph{Random seeds of the analyses}
The draws of the analyses have their own random seeds, fixed in each script: 2018 for the draws of E1 and B1
(\texttt{np.random.default\_rng(2018)}), 2026 for E5 and for every rliable bootstrap
(\cref{sec:reproductibilite-bootstrap}). Two executions of an analysis give a file identical byte for byte.

\paragraph{Each script names its random seeds}
The analysis scripts first read every file present. With random seeds 10 to 19 from deviation~12, E6 would have
gone from ten to twenty random seeds without saying so. Since deviation~14, item~f, each script names its
random seeds and stops if one is missing. The published results did not move.

\section{Evaluation, held out}\label{sec:evaluation}

\paragraph{The evaluation protocol}
After each training block, the \texttt{evaluate} function of \file{sim/run.py} plays ten episodes. For episode
$i$, from 0 to 9:
\begin{enumerate}
  \item a fresh environment is created (\texttt{gym.make(TASK)}), separate from the training one;
  \item it is reset with random seed $\num{10000}+i$;
  \item the policy chooses each action without randomness (\texttt{deterministic=True}): the mean action for
    PPO, SAC and CrossQ, instead of a draw from the policy's distribution;
  \item the episode runs until the body falls or until \num{1000}~steps, then its return is recorded.
\end{enumerate}
The function returns the mean and the standard deviation of the ten returns (\texttt{np.std}, without Bessel's
correction). For the programs tuned like the Zoo, each observation goes through the training normalisation,
whose statistics stay frozen during evaluation. A deterministic policy evaluated on ten fixed random seeds
always gives the same score. The variation among the ten episodes comes from the starting positions alone.

\paragraph{The score of a run}
The score of a run is the mean of the ten returns at the last evaluation, at its program's budget: the
\texttt{eval\_mean} column of the last line of its file. The protocol fixed it as the main measure. Experiment
B3 tries two others.

\paragraph{A modified physics, at the end}
At the end of training, the final policy plays ten more episodes, on the same evaluation random seeds, in a
modified physics. The code modifies the environment's MuJoCo model before the reset:
\begin{itemize}
  \item the mass of every body is multiplied by 1.2 (\texttt{body\_mass[:] *= 1.2});
  \item the sliding friction coefficient of every geometry, the floor included, is multiplied by 0.8
    (\texttt{geom\_friction[:, 0] *= 0.8}).
\end{itemize}
The mean of the ten returns goes into the \texttt{perturbed\_mean} column of the last line. Experiment E7
uses it.

\paragraph{The physics modified by degrees}
Deviation~12, item~ii, re-evaluates each final policy in other physics, with no training at all
(\file{sim/dose.py}). One quantity changes at a time:
\begin{itemize}
  \item the masses, multiplied by 0.7, 0.8, 0.9, 1.0, 1.1, 1.2 and 1.3, from 30\% lighter to 30\% heavier;
  \item the sliding friction, multiplied by 0.6, 0.7, 0.8, 0.9, 1.1 and 1.2, from 40\% weaker to 20\%
    stronger.
\end{itemize}
Each physics is played on the same ten evaluation random seeds, with the same evaluation code. The factor 1.0
serves as a control: a correct re-evaluation gives back there the final score written by \file{sim/run.py}. The
result of each policy is a file
\file{results/dose/}\allowbreak\file{<program>/}\allowbreak\file{<task>/}\allowbreak\file{seed<k>.json}.
The PPO policies of the main grid have no observation normalisation: \file{sim/dose.py} evaluates them as
they are.

\section{The run format}\label{sec:format}

\paragraph{One CSV file per run}
Each run writes a file
\file{results/}\allowbreak\file{<program>/}\allowbreak\file{<task>/}\allowbreak\file{seed<k>.csv}, rewritten
after each evaluation. Its columns are
\texttt{step,}\allowbreak\texttt{eval\_mean,}\allowbreak\texttt{eval\_std,}\allowbreak\texttt{perturbed\_mean,}\allowbreak\texttt{wall}:
\begin{itemize}
  \item \texttt{step}, the number of steps announced at this evaluation, in increments of \num{50000};
  \item \texttt{eval\_mean} and \texttt{eval\_std}, the mean and the standard deviation of the returns of the
    ten evaluation episodes;
  \item \texttt{perturbed\_mean}, the mean of the returns in the modified physics, empty except on the last
    line;
  \item \texttt{wall}, the seconds elapsed since the start of the run.
\end{itemize}
A PPO or TRPO curve has twenty lines, a SAC or CrossQ curve has four. The module
\file{src/pineau\_rl/runs.py} reads this format. Its function \texttt{read\_run} checks the required columns
and that the steps strictly increase. It treats as unfinished a run whose last line is not at the program's
budget. On request, it also cuts a curve at a given step (the equal-budget reading). The function
\texttt{write\_run} writes this format from any library.

\paragraph{The final policies}
Each run saves its final policy in
\file{models/}\allowbreak\file{<program>/}\allowbreak\file{<task>/}\allowbreak\file{seed<k>.zip}, with the
normalisation statistics alongside for the programs tuned like the Zoo (\file{.vecnorm.pkl}). The file
\file{sim/models/MANIFEST.sha256} gives the SHA-256 fingerprint of 250~policies, taken on 15~November~2026 on
the machine that trained them. The PPO, TRPO and E6 policies are in the repository. Those of SAC and CrossQ,
too heavy for git, are attached to the \texttt{policies-v1} release of the repository on GitHub. The script
\file{sim/fetch\_models.py} downloads them one at a time, separated by a random pause of 1.5 to 4~seconds, then
checks each archive and each policy against the manifest. The option \texttt{-{}-check} only checks the
policies already present.

\paragraph{The logs}
Each launch writes a log. Since deviation~14, each attempt keeps its own
(\file{logs/<task>.attempt<n>.log}, never overwritten) and appends a line to the ledger
\file{logs/attempts.jsonl}: task, attempt number, start, end, exit code, SHA-256 fingerprint of the result. For
the earlier period, the ledger \file{logs/attempts\_reconstructed.csv} rebuilds one line per attempt from the
dates of the logs, the outputs of the launchers and the git history (\cref{sec:calcul}).

\section{The statistics}\label{sec:statistiques}

This section defines each statistic used, in the order in which \cref{chap:06-experiences-resultats} needs
them. The functions are in \file{src/pineau\_rl/stats.py} and \file{src/pineau\_rl/}\allowbreak\file{bootstrap.py}.

\subsection{The interquartile mean}\label{sec:iqm}
The interquartile mean (IQM) discards the lowest quarter and the highest quarter of the scores, then averages
the rest \parencite{agarwal2021precipice}. It resists an outlying run better than a mean does and is more
precise than a median. rliable computes it with SciPy's trimmed mean at 25\% on each side
(\texttt{trim\_mean(x, 0.25)}), which removes $\lfloor n/4\rfloor$ scores at each end. With the sorted scores
$x_{(1)}\le\dots\le x_{(n)}$ and $g=\lfloor n/4\rfloor$:
\begin{equation}\label{eq:iqm}
  \mathrm{IQM}(x) = \frac{1}{n-2g}\sum_{i=g+1}^{n-g} x_{(i)} .
\end{equation}
With ten random seeds, $g=2$: the IQM is the mean of the six middle scores. With twenty, $g=5$ and ten scores
remain. With thirty, $g=7$ and sixteen remain.

\subsection{The stratified bootstrap interval}\label{sec:bootstrap}
The uncertainty of an IQM is measured by the bootstrap \parencite{efron1979bootstrap}. We draw $n$ scores with
replacement among the $n$ observed, recompute the IQM, then repeat $B$ times. The 95\% interval is formed by
the 2.5th and 97.5th percentiles of the $B$ values obtained (the percentile method, rliable's default):
\begin{equation}\label{eq:bootstrap}
  \bigl[\,q_{0.025}\bigl(\mathrm{IQM}(x^{*1}),\dots,\mathrm{IQM}(x^{*B})\bigr),\;
  q_{0.975}\bigl(\mathrm{IQM}(x^{*1}),\dots,\mathrm{IQM}(x^{*B})\bigr)\,\bigr],
  \qquad B = \num{2000}.
\end{equation}
rliable calls this bootstrap \enquote{stratified}: with several tasks, it resamples the runs within each
task, without mixing tasks. Each of our calls handles one task, where stratification changes nothing:
the bootstrap amounts to resampling the runs of one program.

\subsection{The probability of improvement}\label{sec:poi}
The probability of improvement of a program $X$ over a program $Y$ is the probability that a run of $X$, drawn
at random, beats a run of $Y$, drawn at random, ties counting as one half \parencite{agarwal2021precipice}.
With $n_x$ scores $x_i$ and $n_y$ scores $y_j$:
\begin{equation}\label{eq:poi}
  P(X>Y) = \frac{1}{n_x n_y}\sum_{i=1}^{n_x}\sum_{j=1}^{n_y}
  \Bigl(\mathbf{1}[x_i>y_j] + \tfrac12\,\mathbf{1}[x_i=y_j]\Bigr) = \frac{U}{n_x n_y},
\end{equation}
where $U$ is the Mann-Whitney statistic \parencite{mann1947test}. rliable computes it this way
(\texttt{scipy.stats.mannwhitneyu}). It equals 0.5 when the two programs are even. Its 95\% interval comes from
the same bootstrap as the IQM: the runs of $X$ and those of $Y$ are resampled separately, $B=\num{2000}$ times,
then the 2.5th and 97.5th percentiles are taken.

\subsection{The winner rule}\label{sec:regle}
The protocol fixes one rule for declaring that one program beats another. We call it \emph{the
settled-winner rule}. A gap is settled when one of two conditions holds:
\begin{itemize}
  \item the 95\% intervals of the two IQMs do not overlap;
  \item the 95\% interval of the probability of improvement excludes 0.5.
\end{itemize}
Otherwise no winner is declared, whatever the gap between the IQMs. E3 applies both conditions
(\file{e3.py}). E2, E2 with Bonferroni and B3 use the second: a pair is called \enquote{settled} there when the
interval of its probability of improvement excludes 0.5. E7 and the physics by degrees use the analogue for a
ratio: a loss is settled when the interval of the IQM of the ratio stays below 1.

\subsection{Student's and Welch's t-tests}\label{sec:welch}
The t-test compares two means. For two groups of sizes $n_a$ and $n_b$, with means $\bar a$ and $\bar b$ and
sample variances $s_a^2$ and $s_b^2$ (with Bessel's correction), Student's test
\parencite{student1908probable} assumes equal variances and computes
\begin{equation}\label{eq:student}
  t = \frac{\bar a-\bar b}{s_p\sqrt{1/n_a+1/n_b}},\qquad
  s_p^2=\frac{(n_a-1)s_a^2+(n_b-1)s_b^2}{n_a+n_b-2},
\end{equation}
with $n_a+n_b-2$ degrees of freedom. Welch's test \parencite{welch1947generalization} does not assume equal
variances:
\begin{equation}\label{eq:welch}
  t = \frac{\bar a-\bar b}{\sqrt{s_a^2/n_a+s_b^2/n_b}},\qquad
  \nu = \frac{\bigl(s_a^2/n_a+s_b^2/n_b\bigr)^2}{\dfrac{(s_a^2/n_a)^2}{n_a-1}+\dfrac{(s_b^2/n_b)^2}{n_b-1}} .
\end{equation}
The two-sided $p$-value is $2\,\bigl(1-F_\nu(|t|)\bigr)$, where $F_\nu$ is the cumulative distribution function
of Student's distribution with $\nu$ degrees of freedom. Two groups of five give $\nu\ge 4$, since $\nu$ is at
least the smaller of the two values $n-1$. All our t-tests are Welch tests (\texttt{scipy.stats.ttest\_ind},
option \texttt{equal\_var=False}). Experiment E1 puts this test in question, as the protocol planned.

\subsection{Three readings of the 2018 average t}\label{sec:t-moyen}
The caption of figure~5 of \textcite{henderson2018matters} speaks of the \enquote{average 2-sample t-test
across entire training distribution}. It compares two groups of whole learning curves. Let $C^a_{s,k}$ be the
evaluation score of run $s$ of group $a$ at evaluation point $k$, for $k=1,\dots,K$ ($K=20$ for PPO and TRPO).
Our analyses compute three readings of this caption (\texttt{welch\_curve} and \texttt{welch\_pooled}).
\begin{itemize}
  \item \emph{The average t.} A Welch test at each point, on the five scores of each group, gives $t_k$ and
    $\nu_k$. We report $\bar t = \frac1K\sum_k t_k$, converted to $p$ with the mean of the $\nu_k$.
  \item \emph{The average p.} The same computation, summarised by $\bar p = \frac1K\sum_k p_k$.
  \item \emph{The pooled points.} The $5\times K$ points of each group are taken as independent observations,
    then compared by one Welch test, with about 198~degrees of freedom.
\end{itemize}

\subsection{The statistic of Henderson's script}\label{sec:script-henderson}
On 17~June~2017, Peter Henderson published the script \file{from\_same\_dist.py}
\parencite{henderson2017samedist}, a companion to \textcite{islam2017reproducibility}. The script computes
something other than the three readings above. The function \texttt{henderson\_script} reproduces it on our
files. Let $M_{s,k}$ be the \texttt{eval\_mean} column and $S_{s,k}$ the \texttt{eval\_std} column of run $s$
at point $k$. Let $A$ and $B$ be the two groups of five runs.
\begin{enumerate}
  \item The script averages the curves of each group, column by column:
    $\bar M^A_k=\frac15\sum_{s\in A}M_{s,k}$ and $\bar S^A_k=\frac15\sum_{s\in A}S_{s,k}$.
  \item At each point $k$, it applies a Welch test built from these statistics alone
    (\texttt{scipy.stats.ttest\_ind\_from\_stats}), with $n$ episodes per group:
    \begin{equation}\label{eq:script}
      t_k=\frac{\bar M^A_k-\bar M^B_k}{\sqrt{\bigl((\bar S^A_k)^2+(\bar S^B_k)^2\bigr)/n}} .
    \end{equation}
  \item It reports the mean of the $t_k$ and the mean of the $p_k$.
\end{enumerate}
The denominator of \cref{eq:script} contains the standard deviation \emph{between the episodes} of one run. The
variability \emph{between random seeds}, the one the 2018 figure was meant to measure, never enters it. The
five random seeds of a group merge into one mean curve. For our runs, $n=10$, the number of evaluation
episodes. A sensitivity analysis takes $n=5$, the number of trajectories per iteration in the 2018 runs
(\num{5000}~steps per iteration, episodes of \num{1000}~steps). Our $S_{s,k}$ is a standard deviation without
Bessel's correction, passed to the test as it is, just as the script passes its own columns.

\subsection{Corrections for multiple comparisons}\label{sec:multiples}
When $m$ tests are run, the probability that at least one is wrong exceeds the threshold of each. Two
corrections are used.
\begin{itemize}
  \item \emph{Bonferroni} \parencite{dunn1961multiple} raises each interval to the level $1-0.05/m$. The $m$
    intervals then hold jointly at a level of at least 95\%. E2 counts $m=24$ comparisons (6~pairs on
    4~tasks): each interval moves to the level $1-0.05/24$, that is 99.79\% (deviation~18).
  \item \emph{Holm} \parencite{holm1979simple} adjusts sorted $p$-values $p_{(1)}\le\dots\le p_{(m)}$ into
    $\tilde p_{(i)}=\min\bigl(1,\max_{j\le i}(m-j+1)\,p_{(j)}\bigr)$. E6 uses it over its two or three
    comparisons.
\end{itemize}

\subsection{Other tests}\label{sec:autres-tests}
The Shapiro-Wilk test \parencite{shapiro1965analysis} measures how far a distribution departs from the normal
distribution. E1 applies it to the thirty scores of PPO PyTorch. The Mann-Whitney test \parencite{mann1947test}
compares two groups by their ranks. The E6 diagnosis reports it next to the Welch test.

\subsection{Reproducible intervals}\label{sec:reproductibilite-bootstrap}
rliable~1.2.0 accepts a \texttt{random\_state} argument but draws its resamples from NumPy's global generator
(\texttt{np.random.choice}). Six scripts did not give their intervals back identically from one execution to
the next (deviation~13). Each script now seeds this generator at the start (\texttt{np.random.seed(2026)}),
then calls rliable in a fixed order. Every call consumes the global generator, so the interval of a comparison
depends on the calls that precede it in the same script. The correction changed no point estimate and no
verdict. The bounds moved at the level of Monte Carlo noise, by at most 16~points on an IQM of \num{2423} for
E6.

\paragraph{A variant for the E6 diagnosis}
The scripts \file{e6\_diag.py} and \file{e6\_official.py} compute the intervals of the IQMs with
$B=\num{5000}$ resamples. Their probability of improvement has a bootstrap of its own
(\texttt{probability\_of\_improvement\_bootstrap}): \num{5000} draws, the two groups resampled one after the
other with a generator of random seed 2018, without rliable. The same score can therefore receive two slightly
different intervals depending on the script that computes it: the IQM of default PPO PyTorch on
\task{Hopper-v5} has one interval in \file{e6.json} and another in \file{e6\_diag.json}. The report cites the
diagnosis interval, as \file{docs/conclusions.md} does.

\section{The experiments, one by one}\label{sec:plans}

\Cref{tab:experiences} summarises the designs. It holds design facts only: the question, the design, the
number of runs and the moment the analysis was written.

{\small
\begin{longtable}{@{}>{\raggedright\arraybackslash}p{1.3cm}>{\raggedright\arraybackslash}p{3.6cm}
  >{\raggedright\arraybackslash}p{4.45cm}>{\raggedright\arraybackslash}p{1.9cm}
  >{\raggedright\arraybackslash}p{2.6cm}@{}}
\caption{The experiments: the question asked, the design, the runs used and the moment the analysis was
fixed. \enquote{Preregistered} means written before any result of the experiment. Every number in this table
is a design choice.}\label{tab:experiences}\\
\toprule
Experiment & Question & Design & Runs & Analysis fixed \\
\midrule
\endfirsthead
\toprule
Experiment & Question & Design & Runs & Analysis fixed \\
\midrule
\endhead
\bottomrule
\endfoot
E1 & Do two groups of runs of the same program pass for two algorithms? &
PPO PyTorch on \task{HalfCheetah}, 30 random seeds. \num{10000} draws of two groups, Welch test on the final
scores, at random then sorted. & 30 training runs & preregistered (8~November) \\
E1, second statistic & Is the average t on the 2018 curve recovered? &
The same 30 curves, Welch test at each of the 20 points, average t and average p. & the 30 of E1 &
after the E1 results (deviation~3) \\
B1, TRPO & Does the spread of TRPO explain the t of 9? &
TRPO on \task{HalfCheetah}, 30 random seeds, three readings of the average t. & 30 training runs &
before any TRPO training (deviation~4) \\
B1, script & Does Henderson's computation explain the t of 9? &
The 2017 script applied to the TRPO and PPO curves, same draws. & the 30 of E1 and the 30 of TRPO &
after B1 TRPO, before this computation (deviation~6) \\
E2 & Does the leading program change from task to task? &
4 programs on 4 tasks, IQM and interval, probability of improvement per pair. &
160 training runs & details fixed before any CrossQ score (deviations~7 and 14) \\
E2, Bonferroni & Does the E2 verdict survive 24 comparisons? &
The 24 intervals of E2 at a simultaneous level of 95\%. & those of E2 &
after E2, before this computation (deviation~18) \\
E3 & Do two programs of the same algorithm give the same result? &
PPO PyTorch against PPO JAX, 4 tasks, same random seeds. & 80 of those of E2 & preregistered \\
E4 & What happens to the mean when the best 3 runs out of 10 are kept? &
Mean and standard deviation of the 10, then of the best 3, reversals among 24 comparisons. & those of E2 &
before any CrossQ score (deviation~7) \\
E5 & How many random seeds recover the correct ranking? &
For each comparison, \num{10000} draws with replacement of $n$ runs, $n$ from 3 to 20. &
those of E2, the 30 of E1 & before any CrossQ score (deviation~7) \\
E6 & What does a baseline tuned with the published settings change? &
PPO PyTorch on \task{Hopper}, default values against Zoo settings, then diagnosis and official script. &
40 training runs beyond E2, plus 10 official & preregistered, diagnosis after the first result
(deviations~2 and 5) \\
E7 & What does a modified physics cost? &
Ratio of the score in modified physics to the score in training physics, per run. & those of E2 &
preregistered, descriptive ranking added afterwards (deviation~11) \\
Physics by degrees & Does the E7 loss come from one modification alone? &
Masses from 0.7 to 1.3 times, friction from 0.6 to 1.2 times, one quantity at a time. &
320 re-evaluations (with twenty random seeds) & rule written before any computation (deviations~12 and 14) \\
B3 & Does the chosen measure change the ranking? &
Final score, best score, mean of the curve, reversals between measures. & those of E2 &
before any CrossQ score (deviation~8) \\
Twenty random seeds & Do the results at ten random seeds hold? &
Random seeds 10 to 19, same rules, E2, E3, E5, E7, B3, physics by degrees. & 150 more training runs &
before any new training (deviations~12 and 20) \\
Equal budget & Does the E2 ranking depend on PPO's budget? &
PPO read at \num{200000}~steps, E2 and B3 recomputed. & those of E2 &
before any look at these scores (deviation~12) \\
E4 replicated & Does the E4 inflation reproduce? & E4 unchanged on random seeds 10 to 19. &
160 runs of random seeds 10 to 19 & before any new training (deviation~12) \\
\end{longtable}
}

\subsection{E1: same code, different random seed}\label{sec:plan-e1}
E1 trains PPO PyTorch thirty times on \task{HalfCheetah-v5}, one million steps each, random seeds 0 to 29,
changing nothing else. The script \file{sim/analysis/e1.py} reads the thirty final scores, then computes, with
a generator of random seed 2018:
\begin{itemize}
  \item the statistics of the distribution: mean, standard deviation, median, minimum, maximum, coefficient of
    variation (standard deviation divided by the mean), Shapiro-Wilk test;
  \item for $k=3$, 5 and 10, \num{10000} draws of two disjoint groups of $k$ runs, obtained by permuting the
    thirty, then taking the first $k$ against the next $k$. Each draw receives a Welch test and the relative
    gap between the means $|\bar a-\bar b|/\min(\bar a,\bar b)$. The script reports the share of draws where
    $p<0.05$, along with the median and the 95th percentile of the relative gap;
  \item the \enquote{sorted} split: \num{10000} draws of ten runs without replacement, sorted by score, then a
    Welch test of the best five against the worst five. This split shows what sorting the runs before the
    comparison would produce;
  \item the selection of the best three out of ten: over \num{10000} draws of ten runs, the mean of the best
    three divided by the mean of the ten.
\end{itemize}
Since the two groups always come from the same program, every declared difference is a false positive. With
groups drawn at random and a well-calibrated test, the expected share is 5\%.

\paragraph{The second statistic}
\textcite{henderson2018matters} test curves, where we tested final scores (\cref{sec:t-moyen}).
\file{sim/analysis/e1\_curve.py} redoes their statistic on the thirty curves: \num{10000} draws of ten runs,
split into two groups of five at random or sorted by final score, then a Welch test at each of the twenty
evaluation points. We report the share of draws declared different by the average t, by the average p and by
the test at the last point, the 95th percentile of $|\bar t|$ and the share of draws where
$|\bar t|\ge 9.09$. This statistic was added on 9~November, after the results of \file{e1.py} (deviation~3).
It complements the preregistered analysis without replacing it.

\subsection{B1: the puzzle of the t of 9}\label{sec:plan-b1}
For TRPO on \task{HalfCheetah}, \textcite{henderson2018matters} publish $t=-9.09$ and $p=0.0016$ between two
groups of five runs of the same code. Pineau specifies that these groups are not sorted. B1 looks for the
origin of this number in two stages, each written before its computation.

\paragraph{First stage: their algorithm}
Deviation~4, written on 9~November before any TRPO training, launches thirty TRPO runs on
\task{HalfCheetah-v5}, random seeds 0 to 29, one million steps, with the default values of sb3-contrib~2.9.2
and the same evaluation as E1. \file{sim/analysis/b1\_trpo.py} applies the three readings of
\cref{sec:t-moyen} to the TRPO and PPO curves, with the same draws as \file{e1\_curve.py}. The reading rule
was fixed in advance. If the 95th percentile of $|\bar t|$ under random splits comes close to 9 with TRPO, the
spread of the algorithm explains their number. If only the pooled reading comes close, their computation
explains it. If no reading comes close, neither lead explains it: the report then says so.

\paragraph{Second stage: their script}
Deviation~6, written on 9~November at 23:55 after the results of \file{b1\_trpo.py} and before any computation
with the script, notes two facts. The first is arithmetic. The published pair $(t,p)=(-9.09,\,0.0016)$ implies
about 3.4~degrees of freedom in a two-sided test. It rules out the pooled-points reading, with about
198~degrees of freedom, which would give a $p$ of the order of $10^{-16}$. It also rules out the average t
converted to $p$, since a Welch test on five against five has at least 4~degrees of freedom. Only the average
of the $p$ values remains compatible. The second fact is the 2017 script (\cref{sec:script-henderson}).
\file{sim/analysis/b1\_script.py} applies this computation to the thirty TRPO curves and the thirty PPO
curves, with the same draws (generator of random seed 2018, \num{10000} draws of ten runs split five against
five, at random then sorted), for $n=10$ then $n=5$. It also reports, per program, the ratio of the mean
standard deviation between episodes at the last point to the standard deviation of the final scores between
random seeds.

\paragraph{Differences stated in advance}
Our evaluation episodes follow a deterministic policy, whereas the 2017 progress files record the training
episodes, played by the stochastic policy. We have twenty evaluation points, while their curves have one line
per iteration. We cannot prove that this script produced their figure~5. Its name (\enquote{from the same
distribution}) and its output (\enquote{t=\dots, p=\dots}) make it likely.

\subsection{E2: does the ranking change from task to task?}\label{sec:plan-e2}
E2 trains the four programs on the four tasks, ten random seeds each. \file{sim/analysis/e2.py} reads the final
score of each run, then, for each task, computes:
\begin{itemize}
  \item the IQM of each program with its 95\% bootstrap interval (\cref{eq:iqm,eq:bootstrap});
  \item the ranking by IQM;
  \item for each of the six pairs, the probability of improvement of the higher-ranked program over the
    other, with its interval (\cref{eq:poi}) and the flag \enquote{settled} if the interval excludes 0.5;
  \item a performance profile: for 25 thresholds $\tau$ from 0 to 1.2, the share of runs whose score, divided
    by the best IQM of the task, exceeds $\tau$.
\end{itemize}
The \emph{leader} of a task is the program with the highest IQM. The ranking \emph{changes} if the leader is
not the same on every task. The change is \emph{settled} when the leader of one task is beaten, in a
settled way, by another program on another task. The script counts these settled changes: each triple
(task where $X$ leads, task where $X$ loses, program that beats $X$) counts as one. These details were written
on 10~November, after the PPO and SAC results and before looking at any CrossQ result (deviation~7, read with
the correction of deviation~14, item~b).

\paragraph{With the Bonferroni correction}
\file{sim/analysis/e2\_bonferroni.py} makes the same calls, in the same order, after the same random seed: the
bootstrap resamples are identical to those of \file{e2.py}. Only the percentile level changes, to
99.79\%. It counts the settled changes with the same definition (deviation~18).

\subsection{E3: two programs of the same algorithm}\label{sec:plan-e3}
\file{sim/analysis/e3.py} pits PPO PyTorch against PPO JAX on each task, with the same random seeds and the
scores at one million steps. It reports the two IQMs with their intervals, the ratio of the IQMs, the
probability that a PPO JAX run beats a PPO PyTorch run with its interval, the overlap of the intervals, then
the winner under the rule of \cref{sec:regle}.

\subsection{E4: the best three out of ten}\label{sec:plan-e4}
For each program and each task, \file{sim/analysis/e4.py} computes the mean and the standard deviation of the
ten runs, then of the best three, finally the inflation $\bar x_{\text{best 3}}/\bar x_{10}-1$. A
comparison is a pair of programs on one task: there are 24. It \emph{reverses} if the sign of the difference in
means changes when each program keeps only its best three runs. The script reports the median inflation, its
range, the median ratio of the standard deviations and the number of reversals.

\subsection{E5: how many random seeds?}\label{sec:plan-e5}
\file{sim/analysis/e5.py} starts from the 24 comparisons. For each, the correct order is the sign of the
difference between the means of all the available runs. The observed distribution of the runs stands in for
the true distribution of the algorithm's scores. For $n=3$, 5, 10 and 20, we draw $n$ runs with replacement
\num{10000} times for each of the two programs, then count the share of draws that recover the correct order:
\begin{equation}\label{eq:e5}
  \hat P_n=\frac{1}{N}\sum_{r=1}^{N}\mathbf{1}\Bigl[\operatorname{sign}\bigl(\bar x^{*(r)}_n-\bar y^{*(r)}_n\bigr)
  =\operatorname{sign}(\bar x-\bar y)\Bigr],\qquad N=\num{10000}.
\end{equation}
For PPO PyTorch on \task{HalfCheetah}, the analysis at ten random seeds draws from the thirty runs of E1. The script
reports the median of $\hat P_n$ over the 24 comparisons, then over the \emph{close} comparisons, those whose
relative gap between means $|\bar x-\bar y|/\max(|\bar x|,|\bar y|)$ stays below 20\%. Deviation~20 fixed a
defect before the computation at twenty random seeds: the script took as the correct order the mean of the first ten
runs, even with twenty. It now takes the mean of all the chosen random seeds. The analysis at ten random seeds came
back identical to the byte.

\subsection{E6: the tuned baseline}\label{sec:plan-e6}
E6 was meant to illustrate flaw F7, the poorly tuned baseline: the gap a \enquote{new method} can claim for
itself when it alone is tuned. The initial design pits, on \task{Hopper-v5}, ten runs of PPO PyTorch with its
default values (those of E2) against ten runs with the Zoo settings (\file{e6.py}).

\paragraph{The diagnosis}
The first result triggered a diagnosis, added after we saw it (deviation~2). \file{e6\_diag.py} compares five
groups of ten runs: default values on v5, Zoo with the discount-factor error on v5, corrected Zoo on v5,
default values on v4, corrected Zoo on v4. Three comparisons are corrected by Holm: the error (corrected Zoo
against Zoo with the error, on v5), the fair test on the Zoo's version (corrected Zoo against default values,
on v4) and the same test on v5. Each comparison reports the Welch $p$, the Mann-Whitney $p$ and the probability
of improvement with its interval.

\paragraph{The official pipeline}
Deviation~5, written after the diagnosis and before any official training, asks whether the plateau seen in
the first result, a score that stops rising after the first training blocks, comes from the published recipe
or from the way we applied it. The script \file{sim/zoo\_official.sh} runs the collection's official script
ten times, \texttt{python -m rl\_zoo3.train -{}-algo ppo -{}-env Hopper-v4 -{}-seed i}, for $i$ from 0 to 9,
in the separate environment \file{sim/env-zoo}, on one thread, with no evaluation during training
(\texttt{-{}-eval-freq -1}). \file{e6\_official.py} then evaluates each official policy with our protocol: ten
deterministic episodes, evaluation random seeds \num{10000} to \num{10009}, observation normalisation loaded
then frozen, no reward normalisation. Two comparisons, fixed in advance and corrected by Holm, pit the official
policies against our corrected Zoo on v4, then against our default values on v4. The reading rule was written
down. If the official runs clearly exceed \num{1000}, our application was at fault. If they plateau too, the
published recipe does worse than the default values in this version of the libraries. The report then says so
without speaking of a baseline badly tuned by us.

\subsection{E7: the modified physics}\label{sec:plan-e7}
For each run $s$, \file{sim/analysis/e7.py} computes the ratio of the two final scores
\begin{equation}\label{eq:e7}
  \rho_s=\frac{\texttt{perturbed\_mean}_s}{\texttt{eval\_mean}_s},
\end{equation}
then, for each program and each task, the IQM of the $\rho_s$ with its bootstrap interval. The verdict is a
\emph{loss} if the interval stays below 1, a \emph{gain} if it stays above, \enquote{no detectable change}
otherwise. The IQM of the ratios differs from the ratio of the IQMs: two rankings of scores do not show the
share of score kept.

\paragraph{The descriptive ranking}
\file{e7\_rank.py} ranks the four programs on each task by the IQM of the score, in the training physics then
in the modified physics. It says whether the leader changes. This complement, decided after the E7 results,
has no test (deviation~11).

\paragraph{The physics by degrees}
\file{sim/analysis/dose.py} reads the 160 re-evaluation files (320 with twenty random seeds). For each program,
task, quantity and factor, it computes the IQM of the ratio to the score at 1.0 times, with its bootstrap
interval, then the same verdict as E7. For each task and each physics, it ranks the four programs by the IQM
of the score, with no test. Finally it checks that the score at 1.0 times recovers the final score of
\file{run.py} and reports the largest relative gap. The rule dates from deviation~12, item~ii, written before
any computation. The script itself was written when 30 of the 160 re-evaluations already existed, one of which
had been displayed: the commit that called it \enquote{written before the results} was wrong (deviation~14,
item~a). The script has no free parameter.

\subsection{B3: does the measure change the ranking?}\label{sec:plan-b3}
\file{sim/analysis/b3.py} computes three measures per run on the recorded curves, with no training at all:
\begin{itemize}
  \item the final score, at the last point;
  \item the best score reached along the curve;
  \item the mean of the curve, the area under the curve divided by its length.
\end{itemize}
For each task and each measure, it computes the IQM of each program and the ranking. A pair of programs
\emph{reverses} between two measures if the order of their IQMs changes. The reversal is \emph{settled}
if, under each of the two measures, the interval of the probability of improvement excludes 0.5, in opposite
directions. PPO curves have twenty points over one million steps, SAC and CrossQ curves four points over
\num{200000}: the mean of the curve does not cover the same budget (deviation~8).

\subsection{The consolidation: twenty random seeds and equal budget}\label{sec:plan-consolidation}
Deviation~12, written on 15~November before any new training, adds three parts with the same code and the
same settings:
\begin{itemize}
  \item random seeds 10 to 19 for the four programs on the four tasks, that is 150 new training runs, since
    PPO PyTorch on \task{HalfCheetah} already had its own;
  \item the physics by degrees of \cref{sec:plan-e7};
  \item the equal-budget reading of \cref{sec:budgets}.
\end{itemize}
The module \file{sim/analysis/runs.py} reads three environment variables, all inactive by default:
\texttt{SEEDS} (for example \texttt{0:20}), \texttt{BUDGET=equal} and \texttt{SUFFIX}, appended to the name of
the output file. Deviation~20, written on 21~November after all training had finished and before any
computation on random seeds 10 to 19, fixes the list of computations, all launched by \texttt{make analysis}:
\begin{itemize}
  \item on the twenty random seeds (suffix \file{\_s20}): E2, E3, E5, E7, B3, the physics by degrees, the
    descriptive ranking of E7 and the Bonferroni correction of E2;
  \item E4 on random seeds 10 to 19 alone (suffix \file{\_r10}), as an independent replication;
  \item E2 and B3 at equal budget, on ten random seeds (\file{\_equal}) then on twenty (\file{\_equal\_s20}).
\end{itemize}
Before deviation~20, the only scores of random seeds 10 to 19 that had been read were the controls of the physics
by degrees at 1.0 times, which compare each file with itself. The last two additions of the first list were not
in deviation~12: they apply rules already written to the twenty random seeds. They are published whatever they
show. The results at ten random seeds remain the main results.

\section{The queue and the compute}\label{sec:calcul}

\paragraph{The machine}
Every training run used the processor of the Mac of \cref{sec:outils}, shared with other projects. A training
run or a re-evaluation occupies one core (one thread). \Cref{tab:calcul} gives, per program and task, the
number of runs, the steps, the median duration and the core-hours.

\gentable{calcul}

\paragraph{The launchers, in the order they were used}
Several scripts launched the runs. Each skips a run whose final policy already exists.
\begin{itemize}
  \item \file{sim/drive.sh} reads a list of jobs (\file{jobs.txt} by default) and launches $P$ of them at a
    time through \texttt{xargs -P}. A run that fails is relaunched once, with a \file{.retry.log} log.
    \file{drive2.sh} is a copy of it for \file{jobs2.txt}, the corrected E6 runs.
  \item \file{sim/relay.sh} switches queues without killing a run: it waits for the running runs to finish,
    then restarts the main queue at one JAX training run (about two cores) and the B1 TRPO queue at two
    training runs, to stay at four cores.
  \item \file{sim/zoo\_official.sh} waits for the TRPO queue to finish, then launches the ten official E6
    training runs on two lanes of five.
  \item \file{sim/queue.sh} reads \file{jobs\_consolidation.txt}, the list of deviation~12: 160
    \texttt{train} lines and 320 \texttt{dose} lines. It reserves each job with an atomic \texttt{mkdir}, so
    two copies run without getting in each other's way. It keeps the ledger \file{logs/attempts.jsonl}. A stop
    by signal, with exit code 128 or above, is not relaunched (deviation~14, item~g).
  \item \file{sim/worker.sh} is \file{queue.sh} with a ceiling: it reads \file{workers.max} before each job
    and exits if too many copies are running (deviation~15).
  \item \file{sim/scale.sh} sets this ceiling every hour. It sums the compute taken by the other programs on
    the machine (\texttt{\%cpu} from \texttt{ps}) and deduces the free cores, minus one as a margin. It grants
    the project at most \texttt{MAX\_CORES} cores, the GRPO run on the processor taking two. It launches new
    \file{worker.sh} copies when the ceiling rises. It launches none without the watchdog. Its decisions are in
    \file{logs/scale.log}.
  \item \file{sim/watchdog.sh}, the watchdog, measures every 20~minutes the free disk space and the growth of
    swap memory since its launch. It stops every launcher and every run if the disk falls below
    \texttt{MIN\_FREE\_GB} gigabytes (10 by default) or if swap has grown by more than 8~GB. It records each
    stopped job in the ledger and releases its reservation.
\end{itemize}

\paragraph{The number of cores over the days}
Until 15~November, the number of simultaneous runs was set by hand: four cores for the project from
9~November, then two copies of \file{queue.sh} for the consolidation. Deviation~15 (15~November) raised the
ceiling to eight cores, set every hour by \file{scale.sh}. Deviation~19 (19~November) brought it down to two
cores for the whole project, the GRPO run on the processor included. The number of parallel jobs changes no
result: each run is determined by its random seed and runs on a single thread.

\paragraph{Stops and relaunches}
The protocol required a crash to be relaunched once and recorded. The reconstructed ledger
\file{logs/attempts\_reconstructed.csv} counts 292~launches for the 250~training runs that precede the
consolidation: 42~attempts were stopped and none crashed. Of the 42 stops, only 2 had been declared
(deviation~14, item~e). The stops have three causes: a general stop on 8~November, probably the one of the
memory incident of the GRPO pilot, a stop requested on 9~November for an overloaded machine, then four stops
to stay at four cores. Each stop hit every running run at once. No run was relaunched or discarded after its
final score was known, nor for any reason linked to its score. For PPO PyTorch and PPO JAX, the relaunches
recover the evaluation points of the stopped attempts identically. For SAC and CrossQ, the stopped curves were
overwritten, which makes this check impossible. The later stops by the watchdog, on 12, 15 and 16~November,
are described in deviations~10, 14 and 16 (\cref{chap:09-ecarts}).

\paragraph{The chain to the documents}
The runs write their CSV and JSON files. The scripts of \file{sim/analysis/} derive the JSON result files from
them. \texttt{make verify} recomputes every analysis and compares the files byte for byte, then
\file{docs/check\_numbers.py} and \file{article/check\_article.py} check every number in the documents against
these files. \Cref{chap:08-logiciel} describes this chain and draws it (\cref{fig:architecture}).

\section{What was written before the results, what was written after}\label{sec:avant-apres}

A protocol written in advance protects against an analysis chosen after seeing the data. It protects only if it
states honestly what was decided afterwards. The decisions that affect an analysis are sorted below, deviation
by deviation.

\paragraph{Written before any result concerned}
\begin{itemize}
  \item the protocol itself (8~November): E1, E2, E3, E4, E5, E6, E7, their statistics, the settled-winner
    rule;
  \item the budget of SAC and CrossQ brought down to \num{200000}~steps (deviation~1), before the first
    training run;
  \item B1 with TRPO and its three readings (deviation~4), before any TRPO training;
  \item the official E6 pipeline and its reading rule (deviation~5), before any official training;
  \item the computation of Henderson's script (deviation~6), before this computation;
  \item the details of E2, E4 and E5 (deviation~7) and the design of B3 (deviation~8), before looking at any
    CrossQ score, when 5 to 7 CrossQ runs out of 40 had finished, instead of the 3 that deviation~7 stated
    (deviation~14, item~b);
  \item the consolidation (deviation~12), before any new training;
  \item the Bonferroni correction (deviation~18), after E2 but before this computation;
  \item the list of computations at twenty random seeds (deviation~20), before any computation on random seeds 10 to 19.
\end{itemize}

\paragraph{Written after seeing a result}
\begin{itemize}
  \item the E6 diagnosis (deviation~2), triggered by the first E6 result;
  \item the second statistic of E1 (deviation~3), after the results of \file{e1.py};
  \item the descriptive ranking of E7 (deviation~11), after the E7 results, with no test;
  \item the random seed of the global generator of the bootstraps (deviation~13), a technical correction with no
    effect on the verdicts;
  \item the seven corrections of deviation~14: the date of the physics-by-degrees script, the counts of
    deviations~7 and~8, the version of Stable-Baselines3, the exact number of steps, the undeclared relaunches,
    the random seeds named in each script, the ledger of attempts.
\end{itemize}
None of these decisions removed a run, changed a setting or changed the settled-winner rule. The deviations
with no effect on the analysis (organisation, interruptions, number of cores, battery) are in
\cref{chap:09-ecarts}.
